\chapter{Related Work}
\label{ch:related-work}

% This chapter provides a comprehensive review of the theoretical foundations that underpin the S\textsuperscript{4}D system, covering web parsing techniques, data modeling approaches, and similarity and clustering methods specifically applicable to web data analysis. The review examines foundational work in web structure analysis, similarity measurement, clustering algorithms, and related applications that inform the development of the S\textsuperscript{4}D system.
This chapter reviews research related to the structural and visual analysis of web pages, highlighting prior work in web parsing, information extraction, and representation learning. The discussion emphasizes how these approaches inform, contrast with, and complement the S\textsuperscript{4}D system. The review is structured into three major areas: (1) structural understanding of web documents, (2) information extraction and representation approaches, and (3) similarity and clustering methods for web analysis. Finally, we compare S\textsuperscript{4}D with prior systems to highlight its unique contributions.

% \section{Web Parsing: DOM, CSS Analysis, and Layout Patterns}
% \label{sec:web-parsing-theory}

% Web parsing encompasses the comprehensive extraction and analysis of web page structural elements, including Document Object Model (DOM) traversal, Cascading Style Sheets (CSS) property analysis, element placement patterns, and XPath-based content extraction. Modern web parsing has evolved significantly since the early days of the World Wide Web, requiring sophisticated approaches to handle complex DOM structures and diverse CSS styling approaches.

% \subsection{Structural Understanding of Web Documents}

% Chen et al. \cite{chen2021websrc} introduced WebSRC, a comprehensive dataset for web-based structural reading comprehension that highlights the importance of understanding both semantic content and structural organization of web pages. Their work demonstrates that effective web page understanding requires systems to comprehend not only textual semantics but also the hierarchical structure of HTML documents. The WebSRC dataset, comprising 400K question-answer pairs from 6.4K web pages with corresponding HTML source code and screenshots, establishes a benchmark for structural web understanding tasks.

% Building on the importance of structural representation, Deng et al. \cite{deng2022domlm} proposed DOM-LM, a novel approach for learning generalizable representations of HTML documents. Their work addresses limitations of existing approaches by encoding both text and DOM tree structure using transformer-based encoders. DOM-LM demonstrates that combining textual content with structural information leads to superior performance in webpage understanding tasks, including attribute extraction, open information extraction, and question answering. The model shows particularly strong generalization capabilities in few-shot and zero-shot settings, making it suitable for real-world applications with limited labeled data.

\section{Structural Understanding of Web Documents}
\label{sec:structural-understanding}

One of the earliest challenges in web analysis has been the ability to combine textual semantics with structural information. Chen et al.~\cite{chen2021websrc} introduced \textbf{WebSRC}, a large-scale dataset for structural reading comprehension. Their work demonstrates that effective web understanding requires not only semantic interpretation but also the hierarchical organization of HTML elements. WebSRC provides 400K question--answer pairs across 6.4K web pages, combining HTML source code, screenshots, and metadata. While WebSRC focuses on question answering, its emphasis on DOM-aware structural comprehension is relevant to S\textsuperscript{4}D’s structural modeling. Unlike WebSRC, however, S\textsuperscript{4}D does not address query answering but instead targets structural and stylistic similarity across domains.

Building on this idea of structural representation, Deng et al.~\cite{deng2022domlm} proposed \textbf{DOM-LM}, a transformer-based model that encodes both textual content and DOM structure to learn generalizable representations of web pages. Their approach addresses the shortcomings of methods that rely solely on visual rendering or plain-text representations. DOM-LM achieves strong generalization in few-shot and zero-shot tasks such as attribute extraction and information retrieval. In contrast, S\textsuperscript{4}D applies structural encoding not for semantic extraction but to identify recurring patterns of element arrangement across domains, enabling inter-domain similarity assessment.


% \subsection{Information Extraction from Web Pages}

% Lockard et al. \cite{lockard2020zeroshotceres} tackled the challenge of zero-shot relation extraction from semi-structured webpages through their ZeroShotCeres system. Their approach recognizes that textual semantics in web documents are often augmented with visual elements including layout, font size, and color. The system employs graph neural networks to build rich representations of text fields and their relationships, enabling generalization to previously unseen website templates. This work is particularly relevant to the S\textsuperscript{4}D system's goal of analyzing structural and visual similarities across different website domains.

% The challenge of web-scale information extraction was addressed by Gulhane et al. \cite{gulhane2011vertex} through their Vertex system. Vertex represents one of the first systems to achieve high-precision information extraction at web scale, processing over 250 million records from more than 200 websites. The system's approach to grouping similar structured pages, learning robust XPath-based extraction rules, and detecting site changes provides valuable insights for large-scale web analysis systems like S\textsuperscript{4}D.
\section{Information Extraction and Representation}
\label{sec:information-extraction}

The extraction of structured information from web pages has been another central research theme. Gulhane et al.~\cite{gulhane2011vertex} developed \textbf{Vertex}, a wrapper induction system deployed at Yahoo! for large-scale record extraction. Vertex introduced algorithms for grouping similar pages, generating XPath-based extraction rules, and adapting to structural changes. It achieved web-scale deployment, extracting over 250 million records from more than 200 websites. Vertex demonstrates the scalability of web information extraction but is tailored to structured record retrieval. By contrast, S\textsuperscript{4}D does not generate wrappers or extract database-like content; it focuses instead on comparing visual and structural layouts across entire domains.

More recently, Lockard et al.~\cite{lockard2020zeroshotceres} proposed \textbf{ZeroShotCeres}, which tackles the challenge of zero-shot relation extraction from semi-structured web pages. Their model employs graph neural networks to capture both textual semantics and visual cues such as layout, font size, and color. This enables generalization to previously unseen templates without the need for template-specific training data. ZeroShotCeres highlights the importance of structural and visual cues for generalizable web analysis. While their focus is relation extraction, S\textsuperscript{4}D instead leverages structural and visual properties to compute domain-level similarity, not entity-level relations.

\section{Similarity and Clustering in Web Analysis}
\label{sec:similarity-clustering}

Similarity measurement and clustering techniques have been applied in web analysis to group documents, identify templates, and extract structure. Traditional approaches often prioritize textual similarity, but recent advances emphasize structural and visual dimensions. For example, graph-based models such as ZeroShotCeres~\cite{lockard2020zeroshotceres} highlight the utility of representing web pages as graphs where nodes correspond to DOM elements and edges represent structural or semantic relationships.

The S\textsuperscript{4}D system extends these ideas by combining multiple similarity dimensions---structural, visual, positional, and semantic---and introducing role vector similarity to capture the functional roles of elements within a document. For clustering, S\textsuperscript{4}D adopts HDBSCAN~\cite{mcinnes2017hdbscan}, which automatically determines the number of clusters, handles noise, and reveals hierarchical patterns. This makes it particularly suited to web data, which is both high-dimensional and heterogeneous. Unlike prior work, S\textsuperscript{4}D applies clustering not at the level of individual pages but across domains, enabling the analysis of cross-domain stylistic and structural similarity.

\section{Comparison with Prior Work}
\label{sec:comparison}

The reviewed works demonstrate significant progress in web analysis, yet they address different objectives from S\textsuperscript{4}D:

\begin{itemize}
    \item \textbf{WebSRC}~\cite{chen2021websrc}: Focuses on question answering over structured web data, highlighting the need for structural comprehension. S\textsuperscript{4}D instead uses structure for cross-domain similarity analysis.
    \item \textbf{DOM-LM}~\cite{deng2022domlm}: Learns generalizable document embeddings from DOM+text representations. S\textsuperscript{4}D leverages DOM+CSS for unsupervised clustering rather than representation learning.
    \item \textbf{Vertex}~\cite{gulhane2011vertex}: Extracts structured records at web scale via wrapper induction. S\textsuperscript{4}D does not extract records but compares layout and style across domains.
    \item \textbf{ZeroShotCeres}~\cite{lockard2020zeroshotceres}: Performs zero-shot relation extraction using structural and visual features. S\textsuperscript{4}D uses similar structural/visual cues but for similarity computation and domain-level clustering.
\end{itemize}

\section{Research Gaps and Opportunities}
\label{sec:research-gaps}

The literature reveals several research gaps that S\textsuperscript{4}D addresses:

\begin{itemize}
    \item \textbf{Cross-domain similarity}: None of the reviewed systems compute similarity matrices across domains, which is central to S\textsuperscript{4}D.
    \item \textbf{Multi-dimensional analysis}: Existing work typically emphasizes content or relations; S\textsuperscript{4}D integrates structural, visual, and positional dimensions.
    \item \textbf{Clustering across stylistic patterns}: Prior systems rarely apply clustering to identify design-level commonalities across domains. S\textsuperscript{4}D introduces a multi-level clustering approach for this purpose.
    \item \textbf{Scalability to modern web design}: Systems like Vertex demonstrate scale for extraction tasks, but scalable cross-domain similarity analysis remains underexplored.
\end{itemize}

% 11/09/2025 mexri edw

% \subsection{DOM Tree Analysis and Processing}

% The Document Object Model (DOM) serves as the foundation for web page structure analysis. Recent advances in DOM-based analysis focus on extracting meaningful patterns from the hierarchical organization of HTML elements. The S\textsuperscript{4}D system builds upon these foundations by implementing recursive DOM traversal algorithms that capture both structural and stylistic properties at multiple depth levels.

% Current approaches to DOM analysis often struggle with the complexity and variability of modern web pages. The S\textsuperscript{4}D system addresses these challenges through its multi-level analysis framework, which processes DOM structures at depths ranging from 1 to 10 levels, allowing for fine-grained structural comparison while maintaining computational efficiency.


\section{Data Modeling for Web}
\label{sec:data-modeling-theory}

Data modeling for web applications involves the systematic representation of web page elements, their attributes, and relationships in structured formats suitable for computational analysis. This includes the transformation of HTML/CSS data into machine-readable formats, the creation of inverted indices for efficient property lookup, and the development of hierarchical data structures that preserve DOM relationships while enabling scalable analysis.

\subsection{Web Data Representation Challenges}

Modern web data modeling faces several key challenges including the heterogeneous nature of HTML structures, the diversity of CSS styling approaches, and the need for normalized representations that enable meaningful comparisons across different websites and design frameworks.

\subsection{JSON Structuring for Web Analysis}

The structured representation of web data through JSON formats provides a foundation for systematic analysis. This includes the creation of hierarchical document models, attribute-value mappings, and metadata integration that supports both individual element analysis and cross-document comparisons.

\section{Similarity and Clustering Methods for Web Data}
\label{sec:similarity-clustering-theory}

Similarity measurement and clustering for web data represent specialized applications of machine learning techniques adapted for the unique characteristics of web content. The challenge lies in developing metrics that can capture the multifaceted nature of web page similarity, including structural, visual, and semantic dimensions while handling the high-dimensional and sparse nature of web feature spaces.

\subsection{Multi-dimensional Similarity Approaches}

Traditional similarity measures focus primarily on textual content, often overlooking the rich structural and visual information present in web pages. The S\textsuperscript{4}D system advances this field by implementing a comprehensive similarity framework that considers multiple dimensions:

\begin{itemize}
    \item \textbf{Structural Similarity}: Based on DOM tree organization and element hierarchies
    \item \textbf{Visual Similarity}: Derived from CSS properties and styling attributes
    \item \textbf{Positional Similarity}: Considering element placement and layout patterns
    \item \textbf{Semantic Similarity}: Incorporating content-based analysis where applicable
\end{itemize}

\subsection{Graph-based Similarity Measures}

Graph neural network approaches, as demonstrated in ZeroShotCeres \cite{lockard2020zeroshotceres}, provide powerful frameworks for capturing complex relationships between web page elements. These approaches represent web pages as graphs where nodes correspond to DOM elements and edges represent structural or semantic relationships.

\subsection{Role Vector Similarity}

The S\textsuperscript{4}D system introduces role vector similarity, which characterizes elements based on their functional roles within the document structure. This approach extends beyond simple structural matching by considering how elements relate to their neighbors and contribute to the overall page organization.

\section{Clustering Algorithms for Web Data}
\label{sec:clustering-algorithms}

Web data presents unique challenges for clustering algorithms due to its high dimensionality, sparsity, and heterogeneous nature. Effective clustering of web pages requires algorithms that can handle these characteristics while producing meaningful and interpretable results.

\subsection{Density-Based Clustering}

The S\textsuperscript{4}D system employs HDBSCAN (Hierarchical Density-Based Spatial Clustering of Applications with Noise) \cite{mcinnes2017hdbscan} as its primary clustering algorithm. HDBSCAN provides several advantages for web data clustering:

\begin{itemize}
    \item \textbf{Automatic cluster number detection}: Eliminates the need for pre-specifying the number of clusters
    \item \textbf{Noise handling}: Effectively identifies and separates outlier web pages
    \item \textbf{Hierarchical structure}: Reveals nested clustering patterns common in web data
    \item \textbf{Parameter robustness}: Less sensitive to parameter selection compared to traditional clustering methods
\end{itemize}

\subsection{Cluster Validation}

The system implements DBCV (Density-Based Cluster Validation) scoring \cite{moulavi2014dbcv} to assess clustering quality. DBCV provides a measure of cluster validity that considers both within-cluster cohesion and between-cluster separation, making it particularly suitable for evaluating the quality of web page groupings.

\subsection{Multi-level Clustering Analysis}

A key innovation in the S\textsuperscript{4}D system is its multi-level clustering approach, which performs clustering analysis at different DOM depth levels. This approach recognizes that different levels of structural granularity may reveal different clustering patterns, allowing for comprehensive analysis of website similarities.

\section{Feature Engineering for Web Analysis}
\label{sec:feature-engineering}

Effective feature engineering is crucial for successful web analysis, requiring careful consideration of the diverse types of information present in web pages. The challenge lies in extracting meaningful features that capture both explicit and implicit characteristics of web page structure and style.

\subsection{CSS Property Standardization}

The S\textsuperscript{4}D system implements comprehensive CSS property processing and normalization to handle the diversity of styling approaches across websites. Key aspects include:

\begin{itemize}
    \item \textbf{Unit normalization}: Converting different measurement units (px, \%, em, rem) to standardized formats
    \item \textbf{Color space standardization}: Processing RGB, RGBA, HSL, and named color values
    \item \textbf{Vendor prefix handling}: Managing browser-specific CSS properties
    \item \textbf{Property filtering}: Focusing on structurally and visually significant properties
\end{itemize}

\subsection{Attribute Selection and Filtering}

Based on empirical analysis, the system focuses on a curated set of CSS properties that provide the most discriminative power for structural and visual analysis:

\begin{itemize}
    \item \textbf{Typography properties}: font-size, font-weight, font-style, text-align
    \item \textbf{Layout properties}: margin, padding, border properties
    \item \textbf{Visual properties}: background-color, color, text-decoration
    \item \textbf{Interaction properties}: cursor, display properties
\end{itemize}

\subsection{Vector Space Representation}

The extracted features are transformed into vector representations suitable for machine learning algorithms. The system employs techniques including:

\begin{itemize}
    \item \textbf{Categorical encoding}: Converting discrete CSS values to numerical representations
    \item \textbf{Dimensionality reduction}: Managing high-dimensional feature spaces
    \item \textbf{Normalization}: Ensuring features contribute appropriately to similarity calculations
\end{itemize}

\section{Similarity Measures for Web Data}
\label{sec:similarity-measures}

Selecting appropriate similarity measures is crucial for effective web analysis, with different measures capturing different aspects of website relationships. The S\textsuperscript{4}D system implements multiple similarity metrics to provide comprehensive analysis capabilities.

\subsection{Cosine Similarity}

Cosine similarity serves as the primary similarity measure in the S\textsuperscript{4}D system, offering several advantages for web data analysis:

\begin{itemize}
    \item \textbf{Magnitude invariance}: Focuses on structural patterns rather than absolute counts
    \item \textbf{High-dimensional efficiency}: Performs well with sparse, high-dimensional feature vectors
    \item \textbf{Interpretability}: Provides intuitive similarity scores between 0 and 1
\end{itemize}

\subsection{Earth Mover's Distance (EMD)}

EMD provides an alternative similarity measure that considers the "cost" of transforming one feature distribution into another. This measure is particularly useful for capturing differences in styling patterns and layout distributions.

\subsection{Hausdorff Distance}

Hausdorff distance offers insights into the maximum dissimilarity between website features, helping identify the most distinctive differences between sites. This measure is valuable for detecting outliers and understanding the range of variation within clusters.

\subsection{Combined Similarity Metrics}

The S\textsuperscript{4}D system combines multiple similarity measures using geometric mean aggregation to create robust similarity assessments that capture different aspects of website relationships.



\section{Applications and Use Cases}
\label{sec:applications}

Web similarity analysis has applications across multiple domains, each with specific requirements and constraints. The diversity of applications highlights the importance of flexible, configurable analysis frameworks like S\textsuperscript{4}D.

\subsection{Web Mining and Content Analysis}

Large-scale web mining operations, such as those implemented in the Vertex system \cite{gulhane2011vertex}, require robust approaches to grouping and analyzing similar web pages. The S\textsuperscript{4}D system's clustering capabilities support similar applications by identifying structurally related pages across different domains.

\subsection{User Interface Research}

Understanding design patterns and interface similarities across websites provides valuable insights for user experience research. The S\textsuperscript{4}D system's focus on visual and structural properties makes it particularly suitable for analyzing design trends and interface patterns.

\subsection{Website Classification and Categorization}

Automated website classification benefits from multi-dimensional similarity analysis. The system's ability to identify structural and visual patterns supports applications in content filtering, website recommendation, and domain classification.

\subsection{Design Pattern Analysis}

The system enables researchers and practitioners to identify common design patterns across different website categories, supporting evidence-based design decisions and interface standardization efforts.

\section{Research Gaps and Opportunities}
\label{sec:research-gaps}

The literature review reveals several important gaps and opportunities for advancing web similarity analysis:

\subsection{Integration of Multiple Analysis Dimensions}

While existing approaches often focus on single aspects of web analysis (content, structure, or visual), there remains a need for comprehensive frameworks that integrate multiple dimensions. The S\textsuperscript{4}D system addresses this gap by combining structural, visual, and positional analysis.

\subsection{Scalability and Performance}

Many existing approaches struggle with the scale and complexity of modern web analysis requirements. Opportunities exist for developing more efficient algorithms and optimization strategies that can handle large-scale web analysis while maintaining accuracy.

\subsection{Dynamic Content Handling}

Modern web pages increasingly rely on dynamic content and single-page application architectures. Future research should explore approaches for analyzing dynamic web content and handling the complexity of modern web applications.

\subsection{Cross-domain Generalization}

Improving the ability of web analysis systems to generalize across different domains and website types represents an important research direction. The zero-shot capabilities demonstrated in ZeroShotCeres \cite{lockard2020zeroshotceres} point toward promising approaches for this challenge.

\section{Literature Review Summary}
\label{sec:literature-summary}

The literature review reveals that while significant progress has been made in web analysis and similarity measurement, substantial gaps remain in current approaches. Key findings include:

\begin{itemize}
    \item \textbf{Structural Understanding}: Recent work like WebSRC \cite{chen2021websrc} and DOM-LM \cite{deng2022domlm} demonstrates the critical importance of combining structural and textual analysis for comprehensive web understanding.
    
    \item \textbf{Scalability Challenges}: Systems like Vertex \cite{gulhane2011vertex} show that web-scale analysis is achievable but requires careful attention to performance optimization and system design.
    
    \item \textbf{Generalization Capabilities}: Approaches like ZeroShotCeres \cite{lockard2020zeroshotceres} highlight the value of developing systems that can generalize to previously unseen website templates and domains.
    
    \item \textbf{Multi-dimensional Analysis}: The literature indicates a trend toward more comprehensive analysis frameworks that consider multiple aspects of web page similarity rather than focusing on single dimensions.
    
    \item \textbf{Feature Engineering}: Effective feature extraction and representation learning remain critical challenges, particularly for handling the diversity and complexity of modern web pages.
\end{itemize}

The S\textsuperscript{4}D system builds upon these insights to provide a comprehensive framework for web similarity analysis that addresses many of the identified gaps while providing practical capabilities for real-world applications. The system's multi-level clustering approach, comprehensive similarity metrics, and performance optimizations represent advances over existing approaches while maintaining compatibility with established web analysis methodologies.

\section{Transition to System Design}
\label{sec:transition-to-system-design}

Having established the theoretical foundations for web parsing, data modeling, and similarity/clustering methods for web data, the following chapter presents the practical implementation of these concepts in the S\textsuperscript{4}D system.

\textbf{What Section 3 Will Cover:}
Chapter~\ref{ch:s4d-system} will detail the complete 8-component S\textsuperscript{4}D system including: (1) system architecture overview, (2) comprehensive technology stack (Selenium, Python, React, FastAPI), (3) multi-level web parsing implementation with DOM analysis across 10 levels, (4) JSON-based data modeling for both pre- and post-clustering phases, (5) data mining techniques for transition matrix construction and role vector extraction, (6) multi-metric similarity analysis combining Cosine, EMD, and Hausdorff distances with best-level selection algorithms, (7) interactive presentation systems with React-based visualization and similarity matrix display, and (8) practical usage guidelines with parameter optimization strategies.

\textbf{What Will Not Be Covered:}
Section 3 will not address theoretical justifications for clustering algorithms (covered in this chapter), alternative web parsing methodologies beyond the implemented approach, or comparative theoretical analysis of similarity measures (discussed above). The focus will be exclusively on the practical implementation, system architecture, and operational aspects of the S\textsuperscript{4}D methodology. 